%============================================================
% BAB 4: RANCANGAN DAN IMPLEMENTASI SISTEM
%   4.1 Standar Perusahaan (jika ada)
%       4.1.x sub-standar
%   4.2 Proyek [Nama Proyek 1]
%       4.2.1 Rancangan Sistem
%   4.3 Proyek [Nama Proyek 2]
%       4.3.1 Rancangan Sistem
%   (hapus section yang tidak relevan jika hanya 1 proyek)
%============================================================
\chapter{RANCANGAN DAN IMPLEMENTASI SISTEM}

Bab ini menjelaskan proses perancangan dan implementasi sistem selama kerja
praktik di \NamaPerusahaan\ dalam periode \PeriodeKP. Pembahasan mencakup
standar perusahaan dalam pengembangan perangkat lunak serta implementasi
proyek-proyek yang dikerjakan.

\section{STANDAR PERUSAHAAN}

Jelaskan standar pengembangan perangkat lunak yang diterapkan di \NamaPerusahaan\
dan wajib diikuti selama pelaksanaan kerja praktik. Standar ini menjadi landasan
dalam proses pengembangan semua proyek.

% Sesuaikan sub-section dengan standar yang berlaku di perusahaan.
\subsection{[Standar 1]}

Jelaskan standar version control atau branching model yang diterapkan.

\subsubsection{[Sub-standar]}

Jelaskan detail branching model yang digunakan.

\subsubsection{[Sub-standar]}

Jelaskan standar pesan commit yang diterapkan.

\subsection{[Standar 2]}

Jelaskan standar penamaan atau konvensi database yang diterapkan.

\subsection{[Standar 3]}

Jelaskan standar struktur direktori atau arsitektur proyek yang digunakan.

\section{PROYEK [NAMA PROYEK 1]}

Jelaskan latar belakang, tujuan, dan deskripsi aplikasi dari proyek pertama
yang dikerjakan selama kerja praktik.

\subsection{Rancangan Sistem}

Jelaskan perancangan sistem dari proyek pertama, meliputi arsitektur sistem,
desain basis data, alur data, dan komponen teknis lainnya. Sertakan diagram
yang relevan.

% Contoh penyisipan gambar arsitektur:
% \begin{figure}[H]
%   \centering
%   \includegraphics[width=0.9\textwidth]{assets/gambar/arsitektur_proyek1.png}
%   \caption{Desain Arsitektur Sistem [Nama Proyek 1]}
%   \label{fig:arsitektur_proyek1}
% \end{figure}

\subsubsection{[Sub-rancangan]}

Jelaskan detail sub-komponen rancangan pertama.

\subsubsection{[Sub-rancangan]}

Jelaskan detail sub-komponen rancangan kedua.

% Hapus section berikut jika hanya mengerjakan 1 proyek.
\section{PROYEK [NAMA PROYEK 2]}

Jelaskan latar belakang, tujuan, dan deskripsi aplikasi dari proyek kedua
yang dikerjakan selama kerja praktik.

\subsection{Rancangan Sistem}

Jelaskan perancangan sistem dari proyek kedua, meliputi arsitektur sistem,
desain basis data, alur data, dan komponen teknis lainnya.

% Contoh penyisipan gambar arsitektur:
% \begin{figure}[H]
%   \centering
%   \includegraphics[width=0.9\textwidth]{assets/gambar/arsitektur_proyek2.png}
%   \caption{Desain Arsitektur Sistem [Nama Proyek 2]}
%   \label{fig:arsitektur_proyek2}
% \end{figure}
